\documentclass[../../main.tex]{subfiles}
\graphicspath{{../../image/}{../../figure/}}
\begin{document}

\section{简单波}\label{section:ch07-简单波}

在古典解的框架内，标量守恒律\eqref{align:ch07-标量守恒律}（特征速度 $\lambda=f'(u)$）取形式
\begin{align}
\partial_t u(x,t)+\lambda(u(x,t))\partial_x u(x,t)=0.
\label{align:ch07-标量对流}
\end{align}

如第 6.1 节已经指出，由\eqref{align:ch07-标量对流}，$u$ 沿特征线保持常值，而这又反过来蕴含每条特征线以常速传播，即它是一条直线。结果表明，一般双曲守恒律组\eqref{align:ch07-典范守恒律}容许具有相同特征的特殊解：

\begin{definition}[$i$-简单波]\label{definition:ch07-i简单波}
双曲守恒律组\eqref{align:ch07-典范守恒律}的一个古典 $C^1$ 解 $U$ 称为\emph{$i$-简单波}，如果 $U$ 沿与它关联的任一 $i$-特征线保持常值。
\end{definition}

于是，一个定义在 $\mathbb{R}^2$ 的开子集上、取值于 $\mathscr{O}$ 的 $C^1$ 函数 $U$ 是 $i$-简单波，如果它在满足\eqref{align:ch07-典范守恒律}的同时还满足
\begin{align}
\partial_t U(x,t)+\lambda_i(U(x,t))\partial_x U(x,t)=0.
\label{align:ch07-简单波方程}
\end{align}
特别地，在 $i$-简单波中，每条 $i$-特征线以常速传播，因而是一条直线。

若 $U$ 是 $i$-简单波，把\eqref{align:ch07-典范守恒律}与\eqref{align:ch07-简单波方程}结合，我们推出
\begin{align}
\begin{cases}
\partial_x U(x,t)=a(x,t)R_i(U(x,t))\\
\partial_t U(x,t)=-a(x,t)\lambda_i(U(x,t))R_i(U(x,t)),
\end{cases}
\label{align:ch07-简单波分解}
\end{align}
其中 $a$ 是标量场。反之，任何满足\eqref{align:ch07-简单波分解}的 $C^1$ 函数 $U$ 必然是 $i$-简单波。

还可以用 Riemann 不变量给出简单波的另一种刻画：

\begin{theorem}\label{theorem:ch07-简单波Riemann判据}
\eqref{align:ch07-典范守恒律}的一个古典 $C^1$ 解 $U$ 是 $i$-简单波，当且仅当每个 $i$-Riemann 不变量在 $U$ 的定义域的每个连通分量上为常数。
\end{theorem}

\begin{proof}
对任一 $i$-Riemann 不变量 $w$，$\partial_x w=\mathrm{D}w\partial_x U$，$\partial_t w=\mathrm{D}w\partial_t U$。若 $U$ 是 $i$-简单波，由\eqref{align:ch07-简单波分解}与\eqref{align:ch07-Riemann不变量定义}，$\partial_x w$ 与 $\partial_t w$ 恒为零，因而 $w$ 在 $U$ 定义域的任一连通分量上为常数。

反之，回忆 $i$-Riemann 不变量的梯度张成 $R_i$ 的正交补，我们推断：当对所有 $i$-Riemann 不变量 $w$，$\partial_x w=\mathrm{D}w\partial_x U$ 恒为零时，$\partial_x U$ 必满足 $\eqref{align:ch07-简单波分解}_1$。把 $\eqref{align:ch07-简单波分解}_1$ 代入\eqref{align:ch07-典范守恒律}，我们推出 $\eqref{align:ch07-简单波分解}_2$ 也成立，即 $U$ 是 $i$-简单波。证明完成。
\end{proof}

任一常值函数 $U=\bar{U}$ 按\hyperref[definition:ch07-i简单波]{定义 7.6.1}都可视为对每个 $i=1,\cdots,n$ 的 $i$-简单波。然而，方便的做法是把这类平凡解称为\emph{常状态}，而把"简单波"一词保留给在其定义域的任一开子集上都不为常数的解。下述命题表明简单波是常状态的自然邻域，它以物理而非数学的术语非正式地陈述。假设与结论的精确含义可从证明中提取。

\begin{theorem}\label{theorem:ch07-弱前沿后简单波}
任一传入常状态的弱前沿以某族 $i$ 的常特征速度传播。进而，该前沿的尾迹必然是一个 $i$-简单波。
\end{theorem}

\begin{proof}
设定如下：假定方程组\eqref{align:ch07-典范守恒律}是严格双曲的。$U$ 是一个古典 Lipschitz 解，在其定义域上除沿一条 $C^1$ 曲线 $x=\chi(t)$ 的图像外是 $C^1$ 的。$U$ 在其定义域位于 $\chi$ 图像一侧（譬如右侧）的任一点处为常值 $\bar{U}$。与此相对，$\partial_x U$ 与 $\partial_t U$ 从 $\chi$ 图像左侧达到非零极限。于是按第 1.6 节的术语，$\chi$ 是以速度 $\dot{\chi}=\mathrm{d}\chi/\mathrm{d}t$ 传播的弱前沿。特别地，\eqref{align:ch01-weak-front-jump}
在此退化为
\begin{align}
[\mathrm{D}F(\bar{U})-\dot{\chi}I][[\partial U/\partial\mathcal{N}]]=0,
\label{align:ch07-弱前沿Lax}
\end{align}
它表明 $\dot{\chi}$ 是常数且等于某个 $i$ 的 $\lambda_i(\bar{U})$。

接下来我们证明：在 $\chi$ 图像左侧且充分靠近处，解 $U$ 是 $i$-简单波。由\hyperref[theorem:ch07-简单波Riemann判据]{定理 7.6.2}，只需证明 $n-1$ 个独立的 $i$-Riemann 不变量（记为 $w_1,\cdots,w_{i-1},w_{i+1},\cdots,w_n$）为常数。

对靠近 $\bar{U}$ 的 $U$，向量组 $\{\mathrm{D}w_1(U),\cdots,\mathrm{D}w_{i-1}(U),\mathrm{D}w_{i+1}(U),\cdots,\mathrm{D}w_n(U)\}$ 张成 $R_i(U)$ 的正交补；左特征向量组 $\{L_1(U),\cdots,L_{i-1}(U),L_{i+1}(U),\cdots,L_n(U)\}$ 同样如此。因而存在非奇异 $(n-1)\times(n-1)$ 矩阵 $B(U)$ 使得
\begin{align}
L_j(U)=\sum_{k\neq i}B_{jk}(U)\mathrm{D}w_k(U),\qquad j=1,\cdots,i-1,i+1,\cdots,n.
\label{align:ch07-Lj用Dw}
\end{align}

从左用 $L_j(U)$ 乘\eqref{align:ch07-典范守恒律}，得到
\begin{align}
L_j(U)\partial_t U+\lambda_j(U)L_j(U)\partial_x U=0,\qquad j=1,\cdots,n.
\label{align:ch07-方程组左投影}
\end{align}
合并\eqref{align:ch07-Lj用Dw}与\eqref{align:ch07-方程组左投影}，我们结论
\begin{align}
\sum_{k\neq i}B_{jk}\partial_t w_k+\sum_{k\neq i}\lambda_j B_{jk}\partial_x w_k=0,\qquad j=1,\cdots,i-1,i+1,\cdots,n.
\label{align:ch07-Riemannw方程组}
\end{align}
我们把\eqref{align:ch07-Riemannw方程组}视为关于 $n-1$ 个未知量 $w_1,\cdots,w_{i-1},w_{i+1},\cdots,w_n$ 的 $n-1$ 个一阶线性非齐次方程组。在此意义下，\eqref{align:ch07-Riemannw方程组}是严格双曲的，其特征速度为 $\lambda_1,\cdots,\lambda_{i-1},\lambda_{i+1},\cdots,\lambda_n$。沿 $\chi$ 的图像，$n-1$ 个 Riemann 不变量是常数，即等于它们在 $\bar{U}$ 处的值：$w_1(\bar{U}),\cdots,w_{i-1}(\bar{U}),w_{i+1}(\bar{U}),\cdots,w_n(\bar{U})$。此外，$\chi$ 的图像对方程组\eqref{align:ch07-Riemannw方程组}是非特征的。因而，线性双曲方程组 Cauchy 问题的标准唯一性定理蕴含：\eqref{align:ch07-Riemannw方程组}只可能有一个与 Cauchy 数据相容的解，即平凡解：
$w_1=w_1(\bar{U}),\cdots,w_{i-1}=w_{i-1}(\bar{U}),w_{i+1}=w_{i+1}(\bar{U}),\cdots,w_n=w_n(\bar{U})$。证明完成。
\end{proof}

在 $i$-简单波 $U$（\eqref{align:ch07-典范守恒律}）定义域的任一点 $(x,t)$，我们令 $\xi(x,t)$ 表示与 $U$ 关联的 $i$-特征线在 $(x,t)$ 处的斜率，即
\begin{align}
\xi(x,t)=\lambda_i(U(x,t)).
\label{align:ch07-简单波斜率}
\end{align}
$\xi$ 在斜率为 $\xi$ 的直线方向上的导数为零，即
\begin{align}
\partial_t\xi+\xi\partial_x\xi=0.
\label{align:ch07-斜率Burgers}
\end{align}
因而 $\xi$ 满足 Burgers 方程\eqref{align:ch04-Burgers}
。

在 $U$ 定义域任一点 $(\bar{x},\bar{t})$ 的邻域内，若 $\partial_x\xi(\bar{x},\bar{t})>0$，我们称该 $i$-简单波是\emph{$i$-稀疏波}（即 $i$-特征线发散）；若 $\partial_x\xi(\bar{x},\bar{t})<0$，称其为\emph{$i$-压缩波}（即 $i$-特征线汇聚）。这一术语源于气体动力学的背景。

由于在 $i$-简单波中 $U$ 沿 $i$-特征线保持常值，在 $\partial_x\xi(\bar{x},\bar{t})\neq 0$ 的任一点 $(\bar{x},\bar{t})$ 的小邻域 $\mathscr{K}$ 内，我们可用单个变量 $\xi$ 标记 $U$，即存在一个函数 $V_i$，定义在区间 $(\bar{\xi}-\varepsilon,\bar{\xi}+\varepsilon)$ 上（其中 $\bar{\xi}=\lambda_i(U(\bar{x},\bar{t}))$），取值于 $\mathscr{O}$，使得
\begin{align}
U(x,t)=V_i(\xi(x,t)),\qquad (x,t)\in\mathscr{K}.
\label{align:ch07-简单波参数化}
\end{align}
进而，由\eqref{align:ch07-简单波分解}与\eqref{align:ch07-简单波斜率}，$V_i$ 满足
\begin{align}
\dot{V}_i(\xi)=b(\xi)R_i(V_i(\xi)),\qquad \xi\in(\bar{\xi}-\varepsilon,\bar{\xi}+\varepsilon),
\label{align:ch07-稀疏波曲线ODE}
\end{align}
\begin{align}
\lambda_i(V_i(\xi))=\xi,\qquad \xi\in(\bar{\xi}-\varepsilon,\bar{\xi}+\varepsilon),
\label{align:ch07-稀疏波速度}
\end{align}
其中 $b$ 是标量函数，圆点表示关于 $\xi$ 的导数。

反之，若 $V_i$ 满足\eqref{align:ch07-稀疏波曲线ODE}与\eqref{align:ch07-稀疏波速度}，且 $\xi$ 是\eqref{align:ch07-斜率Burgers}在区间 $(\bar{\xi}-\varepsilon,\bar{\xi}+\varepsilon)$ 中取值的任一 $C^1$ 解，则 $U=V_i(\xi(x,t))$ 是一个 $i$-简单波。上述考虑激发了下述定义。

\begin{definition}[$i$-稀疏波曲线]\label{definition:ch07-稀疏波曲线}
对双曲守恒律组\eqref{align:ch07-典范守恒律}，状态空间 $\mathbb{R}^n$ 中的一条\emph{$i$-稀疏波曲线}是一条曲线 $U=V_i(\cdot)$，其中函数 $V_i$ 满足\eqref{align:ch07-稀疏波曲线ODE}与\eqref{align:ch07-稀疏波速度}。
\end{definition}

稀疏波曲线将成为在第 IX 章中求解 Riemann 问题的主要工具之一。在真正非线性态的邻域内，这些曲线的构造特别简单：

\begin{theorem}\label{theorem:ch07-稀疏波曲线存在唯一}
假定 $\bar{U}\in\mathscr{O}$ 是双曲守恒律组\eqref{align:ch07-典范守恒律}的 $i$-特征族的真正非线性态。那么存在唯一的、通过 $\bar{U}$ 的 $i$-稀疏波曲线 $V_i$。若 $R_i$ 在 $\bar{U}$ 的邻域上通过
\begin{align}
\mathrm{D}\lambda_i(U)R_i(U)=1,
\label{align:ch07-特征向量规范化}
\end{align}
规范化，并且 $V_i$ 通过 $\tau=\xi-\bar{\xi}$（其中 $\bar{\xi}=\lambda_i(\bar{U})$）重新参数化，则 $V_i$ 是常微分方程
\begin{align}
\dot{V}_i=R_i(V_i)
\label{align:ch07-稀疏波ODE}
\end{align}
在初值 $V_i(0)=\bar{U}$ 下的解。特别地，$V_i$ 是 $C^3$。当需要显示该稀疏波曲线的起点时，将使用更明确的记号 $V_i(\tau;\bar{U})$。
\end{theorem}

\begin{proof}
\eqref{align:ch07-稀疏波ODE}的任一解 $V_i$ 显然满足 $b=1$ 时的\eqref{align:ch07-稀疏波曲线ODE}。在 $\xi=\bar{\xi}$，即 $\tau=0$ 时，$\lambda_i(V_i)=\lambda_i(\bar{U})=\bar{\xi}$。进而，$\dot{\lambda}_i(V_i)=\mathrm{D}\lambda_i(V_i)\dot{V}_i=1$，这由\eqref{align:ch07-稀疏波ODE}与\eqref{align:ch07-特征向量规范化}得到。这确立了\eqref{align:ch07-稀疏波速度}并完成证明。
\end{proof}

相反，当 $i$-特征族是线性退化时，关于 $\xi$ 对\eqref{align:ch07-稀疏波速度}求导并把所得方程与\eqref{align:ch07-稀疏波曲线ODE}结合，导出矛盾 $0=1$。在此情形下，任一 $i$-简单波中的 $i$-特征线必然是平行直线。然而，仍然成立的是：任一 $i$-简单波取某条\eqref{align:ch07-稀疏波ODE}的积分曲线上的值。

受\hyperref[theorem:ch07-简单波Riemann判据]{定理 7.6.2}启发，我们可以用 Riemann 不变量刻画稀疏波曲线：

\begin{theorem}\label{theorem:ch07-稀疏波曲线Riemann刻画}
每个 $i$-Riemann 不变量沿方程组\eqref{align:ch07-典范守恒律}的任一 $i$-稀疏波曲线为常数。反之，若 $\bar{U}$ 是\eqref{align:ch07-典范守恒律}的 $i$-特征族的任一真正非线性态，且 $w_1,\cdots,w_{i-1},w_{i+1},\cdots,w_n$ 是 $\bar{U}$ 某邻域上独立的 $i$-Riemann 不变量，则通过 $\bar{U}$ 的 $i$-稀疏波曲线由方程组 $w_j(U)=w_j(\bar{U})$（$j=1,\cdots,i-1,i+1,\cdots,n$）隐式确定。
\end{theorem}

\begin{proof}
任一 $i$-稀疏波曲线 $V_i$ 满足\eqref{align:ch07-稀疏波曲线ODE}。若 $w$ 是\eqref{align:ch07-典范守恒律}的 $i$-Riemann 不变量，从左乘 $\mathrm{D}w(V_i(\xi))$ 于\eqref{align:ch07-稀疏波曲线ODE}并利用\eqref{align:ch07-Riemann不变量定义}，得到 $\dot{w}(V_i(\xi))=0$，即 $w$ 沿 $V_i$ 保持常数。

现假定 $w_1,\cdots,w_{i-1},w_{i+1},\cdots,w_n$ 是 $i$-Riemann 不变量，使得 $\mathrm{D}w_1,\cdots,\mathrm{D}w_{i-1},\mathrm{D}w_{i+1},\cdots,\mathrm{D}w_n$ 线性无关。于是 $n-1$ 个曲面 $w_j(U)=w_j(\bar{U})$（$j=1,\cdots,i-1,i+1,\cdots,n$）横截相交，构成一条通过 $\bar{U}$ 的 $C^1$ 曲线 $V_i$，以弧长 $s$ 参数化；按\hyperref[definition:ch07-iRiemann不变量]{定义 7.3.1}，其切向量 $V_i'(s)$ 必满足 $V_i'(s)=c(s)R_i(V(s))$，其中 $c$ 是非零标量函数。只要 $V_i$ 是 $i$-特征场的真正非线性态，$\lambda_i'(V_i)=\mathrm{D}\lambda_i V_i'=c\mathrm{D}\lambda_i R_i\neq 0$。于是我们可以找到适当的参数化 $s=s(\xi)$，使 $V_i$ 同时满足\eqref{align:ch07-稀疏波曲线ODE}与\eqref{align:ch07-稀疏波速度}。证明完成。
\end{proof}

作为\hyperref[theorem:ch07-稀疏波曲线Riemann刻画]{定理 7.6.6}的应用，我们推断：当 $\sigma''(\bar{u})\neq 0$ 时，方程组\eqref{align:ch07-p系统}通过点 $(\bar{u},\bar{v})$ 的 1- 与 2-稀疏波曲线按 Riemann 不变量\eqref{align:ch07-p系统Riemann}由方程
\begin{align}
v=\bar{v}\pm\int_{\bar{u}}^u\sqrt{\sigma'(\omega)}\,\mathrm{d}\omega
\label{align:ch07-p系统稀疏波}
\end{align}
确定。

类似地，当 $p_{uu}(\bar{u},\bar{s})\neq 0$ 时，方程组\eqref{align:ch07-热弹性}通过点 $(\bar{u},\bar{v},\bar{s})$ 的 1- 与 3-稀疏波曲线按 Riemann 不变量\eqref{align:ch07-热弹性Riemann}由方程
\begin{align}
v=\bar{v}\pm\int_{\bar{u}}^u\sqrt{-p_\omega(\omega,\bar{s})}\,\mathrm{d}\omega,\qquad s=\bar{s}
\label{align:ch07-热弹性稀疏波}
\end{align}
描述。

当方程组\eqref{align:ch07-典范守恒律}配有 Riemann 不变量坐标系 $(w_1,\cdots,w_n)$ 并采用 $W=(w_1,\cdots,w_n)^\top$（而非 $U$）作为状态变量时，稀疏波曲线取非常简单的形式。事实上，由\hyperref[theorem:ch07-稀疏波曲线Riemann刻画]{定理 7.6.6}，通过点 $\bar{W}=(\bar{w}_1,\cdots,\bar{w}_n)^\top$ 的 $i$-稀疏波曲线是直线 $w_j=\bar{w}_j$（$j\neq i$），平行于 $i$ 轴。

\end{document}
